% !TEX root = ../Main.tex
\chapter{休息时间}

\section{分与分之间的间隔}

\state{25 秒发球规则}{
    \begin{itemize}
        \item 每一分结束后，发球手必须在\textbf{25 秒内}开始下一次发球；
        \item 超时将被裁判\textbf{警告}，多次违反会\textbf{判罚一分}。
    \end{itemize}
}

\section{换边休息}

\state{换边的时机}{
    \begin{itemize}
        \item 每盘内的\textbf{单数局}（即第 1、3、5、\ldots 局）结束后，双方\textbf{交换场地}；
        \item 每盘结束后，如果该盘\textbf{总局数为奇数}也要换边（比如 6--3、7--5）；若为偶数则不换边。
    \end{itemize}
}

\state{换边时的休息时间}{
    \begin{itemize}
        \item \textbf{第 1 局结束后}：不休息，\textbf{直接换边继续比赛}；
        \item \textbf{其他单数局结束后}：休息\textbf{不超过 90 秒}；
        \item \textbf{每盘结束后换边时}：休息\textbf{不超过 120 秒}（即 2 分钟）；
        \item \textbf{抢七决胜局中}：当累计得 \textbf{6 分的倍数}（即 6、12、\ldots）时换边，\textbf{不进行休息}。
    \end{itemize}
}

\section{休息时间汇总表}

\begin{center}
\begin{tabular}{@{} >{\bfseries}l l @{}}
    \toprule
    情景 & 休息时长 \\
    \midrule
    每分之间（发球准备） & $\le$ 25 秒 \\
    第 1 局结束 & 不休息（直接换边） \\
    其他单数局结束（换边） & $\le$ 90 秒 \\
    每盘结束（换边） & $\le$ 120 秒 \\
    抢七中每累计 6 分 & 换边不休息 \\
    \bottomrule
\end{tabular}
\end{center}

\rmkb{
    \textbf{易考点}：
    \begin{itemize}
        \item \textbf{25 秒}发球间隔；
        \item 换边发生在\textbf{每盘的单数局结束}；
        \item \textbf{第 1 局结束不休息}，其他单数局结束休息\textbf{90 秒}；
        \item \textbf{每盘结束}休息\textbf{120 秒}；
        \item \textbf{抢七中 6 分倍数换边不休息}。
    \end{itemize}
}
